% Descripción técnica del software CEP. Compilar con:
%   pdflatex DESCRIPCION_SOFTWARE.tex   (dos pasadas)
\documentclass[11pt,a4paper]{article}
\usepackage[T1]{fontenc}
\usepackage[utf8]{inputenc}
\usepackage{lmodern}
\usepackage[spanish,es-noshorthands,es-tabla,es-nodecimaldot]{babel}
\usepackage[margin=2.4cm]{geometry}
\usepackage{microtype}
\usepackage{amsmath,amssymb}
\usepackage{xcolor}
\usepackage{booktabs,tabularx,longtable,array}
\usepackage{enumitem}
\usepackage{titlesec}
\usepackage{fancyhdr}
\usepackage{needspace}
\usepackage{tikz}
\usetikzlibrary{positioning,arrows.meta}
\usepackage[hidelinks]{hyperref}

\definecolor{primario}{HTML}{0F4C5C}
\definecolor{gristexto}{HTML}{40484C}
\definecolor{fondo}{HTML}{EEF4F5}
\definecolor{borde}{HTML}{BFC8CB}

\titleformat{\section}{\needspace{6\baselineskip}\Large\bfseries\color{primario}}{\thesection.}{0.6em}{}
\titleformat{\subsection}{\needspace{5\baselineskip}\large\bfseries\color{primario}}{\thesubsection}{0.6em}{}
\titleformat{\subsubsection}{\needspace{4\baselineskip}\normalsize\bfseries}{}{0em}{}
\titlespacing*{\section}{0pt}{1.6em}{0.7em}
\titlespacing*{\subsection}{0pt}{1.2em}{0.5em}
\titlespacing*{\subsubsection}{0pt}{1em}{0.3em}

\setlength{\parindent}{0pt}
\setlength{\parskip}{0.55em}
\setlist{itemsep=0.25em,topsep=0.3em,leftmargin=1.4em}
\renewcommand{\arraystretch}{1.25}

\pagestyle{fancy}
\fancyhf{}
\fancyhead[L]{\small\color{gristexto}CEP — Control Estadístico de Procesos}
\fancyhead[R]{\small\color{gristexto}Descripción técnica del software}
\fancyfoot[C]{\small\thepage}
\renewcommand{\headrulewidth}{0.4pt}

\newcommand{\bs}{\textbackslash{}}
\newcommand{\ruta}[1]{\texttt{#1}}
\newcolumntype{L}{>{\raggedright\arraybackslash}X}
\newcolumntype{P}[1]{>{\raggedright\arraybackslash}p{#1}}

\begin{document}

% ---------------------------------------------------------------- portada
\begin{titlepage}
\vspace*{3cm}
{\color{primario}\rule{\linewidth}{1.2pt}}\par\vspace{1.2em}
{\fontsize{30}{36}\selectfont\bfseries\color{primario} CEP\par}
\vspace{0.4em}
{\LARGE Control Estadístico de Procesos\par}
\vspace{1.5em}
{\Large Descripción técnica del software\par}
\vspace{1.2em}
{\color{primario}\rule{\linewidth}{1.2pt}}\par
\vspace{2em}
{\large Versión 1.0.1 \hfill Noviembre de 2026\par}
\vfill
\begin{minipage}{\linewidth}\color{gristexto}
Programa de escritorio para Windows (y aplicación web local) que implementa el control
estadístico de procesos de punta a punta: captura de datos, cartas de control en Fase~I y
Fase~II, validación de supuestos, capacidad, diseño de cartas por ARL/ATS, muestreo de
aceptación y reportes con trazabilidad.

\medskip
Este documento describe qué hace el software y cómo está construido. El uso paso a paso está
en el manual de usuario y la instalación en la guía de instalación.
\end{minipage}
\end{titlepage}

{\setlength{\parskip}{0pt}\tableofcontents}
\newpage

% ---------------------------------------------------------------- 1
\section{Propósito y alcance}

CEP convierte las mediciones o conteos de un proceso productivo en información para decidir:
si el proceso es estable (está bajo control estadístico), si es capaz de cumplir la
especificación, qué fracción del producto se espera fuera de ella y con qué plan de muestreo
conviene vigilarlo.

\begin{tabularx}{\linewidth}{@{}P{3cm}L@{}}
\toprule
\textbf{Aspecto} & \textbf{Alcance}\\
\midrule
Datos & Variables continuas (mediciones) y atributos (unidades defectuosas o defectos por unidad).\\
Fases & Fase~I (estimación de límites con depuración de causas asignables) y Fase~II (monitoreo con límites congelados).\\
Usuarios & Dos perfiles: \textbf{Operario} (registro y lectura del proceso en palabras simples) e \textbf{Ingeniería} (análisis completo).\\
Despliegue & Un solo equipo: servidor local en \texttt{127.0.0.1}, base SQLite, sin internet ni servicios externos.\\
Base teórica & Montgomery, \emph{Introduction to Statistical Quality Control}. Los resultados del núcleo estadístico se validan con sus ejemplos resueltos.\\
\bottomrule
\end{tabularx}

% ---------------------------------------------------------------- 2
\section{Arquitectura}

\subsection{Capas}

\begin{center}
\begin{tikzpicture}[
  caja/.style={draw=borde, fill=fondo, rounded corners=2pt, align=center,
               inner sep=6pt, font=\small},
  flecha/.style={-{Stealth[length=2.2mm]}, draw=primario, thick}]
\node[caja, text width=0.86\linewidth] (web) {%
  \textbf{Interfaz} — ventana propia (WebView2 vía pywebview) o navegador en modo servidor\\
  16 pantallas Jinja2 + JavaScript por pantalla + Plotly.js (todo local)};
\node[caja, text width=0.86\linewidth, below=0.9cm of web] (api) {%
  \textbf{\texttt{app/api}} — FastAPI: 72 operaciones REST en 16 grupos (un router por módulo)\\
  \texttt{servicios.py}: combina base de datos y núcleo, sin fórmulas propias};
\node[caja, text width=0.39\linewidth, below left=0.9cm and -0.43\linewidth of api] (db) {%
  \textbf{\texttt{app/db}} — SQLAlchemy 2\\SQLite (un archivo \texttt{cep.db})\\migraciones y respaldos};
\node[caja, text width=0.39\linewidth, below right=0.9cm and -0.43\linewidth of api] (core) {%
  \textbf{\texttt{app/core}} — núcleo estadístico puro\\listas/arrays $\rightarrow$ \texttt{dict} JSON\\sin dependencias de web ni de BD};
\draw[flecha] (web) -- node[right, font=\footnotesize, color=gristexto]{HTTP/JSON (127.0.0.1)} (api);
\draw[flecha] (api.south -| db.north) -- (db.north);
\draw[flecha] (api.south -| core.north) -- (core.north);
\end{tikzpicture}
\end{center}

\textbf{Regla de dependencias.} \texttt{app/core} no importa FastAPI, SQLAlchemy ni nada de las
otras capas. Cada función recibe datos y devuelve diccionarios serializables, y documenta en su
docstring la fórmula que implementa. Esto permite probar la estadística de forma aislada y
reutilizarla en la API, los reportes y la consola (\texttt{datos\_prueba/demo\_core.py}).

\subsection{Tecnologías}

\begin{tabularx}{\linewidth}{@{}P{3.4cm}L@{}}
\toprule
\textbf{Capa} & \textbf{Tecnología (versión instalada)}\\
\midrule
Lenguaje & Python 3.11\\
Cálculo & NumPy 2.4, SciPy 1.17, pandas 3.1, statsmodels 0.15\\
API y servidor & FastAPI 0.165, Uvicorn 0.54, Pydantic (validación de entradas)\\
Persistencia & SQLAlchemy 2.1 sobre SQLite\\
Interfaz & Jinja2 3.1, JavaScript sin framework, Tailwind CSS (local), Plotly.js 2.35 (local)\\
Reportes & LaTeX (\texttt{pdflatex}, MiKTeX/TeX Live) + Jinja2; respaldo con ReportLab 5.1; Excel con openpyxl 3.1; figuras con Matplotlib 3.11\\
Escritorio & pywebview 6.2 (Microsoft Edge WebView2), PyInstaller 6.22, Inno Setup 6\\
Pruebas & pytest (296 pruebas automáticas)\\
\bottomrule
\end{tabularx}

Todos los recursos de la interfaz (Tailwind, Plotly, tipografías League Spartan y Montserrat,
íconos Material Symbols) se sirven desde la propia aplicación: no hay solicitudes a internet
durante la ejecución.

\subsection{Tamaño del código propio}

\begin{tabularx}{\linewidth}{@{}Lrr@{}}
\toprule
\textbf{Capa} & \textbf{Archivos} & \textbf{Líneas}\\
\midrule
\texttt{app/core} (estadística, reportes, plantilla) & 15 & $\approx$ 4\,000\\
\texttt{app/api} (endpoints y servicios) & 19 & $\approx$ 2\,200\\
\texttt{app/db} (modelos, sesión, semilla, respaldos) & 5 & $\approx$ 400\\
\texttt{app/web} (plantillas HTML y JavaScript) & 39 & $\approx$ 5\,800\\
\texttt{tests} & 24 & $\approx$ 2\,700\\
\bottomrule
\end{tabularx}

% ---------------------------------------------------------------- 3
\section{Modos de uso}

\subsection{Programa de escritorio (distribución a usuarios)}
\begin{itemize}
\item \textbf{Instalador:} \texttt{CEP-Setup-1.0.1.exe} ($\approx$ 69~MB). Instala por usuario,
  sin permisos de administrador ($\approx$ 205~MB instalados), crea accesos directos y un
  desinstalador.
\item \textbf{Ejecución:} \texttt{CEP.exe} levanta el servidor en \texttt{127.0.0.1} en un puerto
  libre (8000--8099) y abre una ventana propia con WebView2. Al cerrar la ventana se detiene el
  servidor. Se impide una segunda instancia (mutex de Windows).
\item \textbf{Datos:} \ruta{C:\bs Users\bs <usuario>\bs CEP\bs} (fuera de OneDrive):
  \texttt{cep.db}, \texttt{reportes\bs}, \texttt{respaldos\bs} y \texttt{registro.log}.
  Desinstalar o actualizar no borra esta carpeta.
\item \textbf{Primer arranque:} base vacía; los datos de demostración se cargan a pedido.
\item \textbf{Respaldo automático} de la base en cada apertura (se conservan los últimos 20).
\end{itemize}

\subsection{Modo servidor (desarrollo y revisión)}
\texttt{run.bat} (Windows) o \texttt{run.sh} (Linux/macOS) crean el entorno virtual, instalan
dependencias, siembran los datos de demostración y abren el navegador en
\texttt{http://127.0.0.1:8000}. La base queda en \texttt{data/cep.db}.

\subsection{API REST}
Toda la funcionalidad está disponible como API JSON con documentación interactiva (OpenAPI) en
\texttt{/docs}. La interfaz no tiene lógica estadística propia: consume esta API.

\subsection{Perfiles de la interfaz}

\begin{tabularx}{\linewidth}{@{}P{2.2cm}P{4.6cm}L@{}}
\toprule
\textbf{Perfil} & \textbf{Pantallas} & \textbf{Enfoque}\\
\midrule
Operario & Registrar muestra, Cómo va el proceso, Alertas, Historial de muestras &
  Lenguaje simple: estado del proceso en palabras («El proceso va bien», «Revisar», «Detener y
  avisar»), qué hacer, una sola gráfica y alertas explicadas sin códigos de reglas.\\
Ingeniería & 12 pantallas (ver sección~\ref{sec:interfaz}) &
  Análisis completo, configuración, diseño, reportes y administración de la base.\\
\bottomrule
\end{tabularx}

No hay inicio de sesión: el nombre del usuario se indica en el encabezado y viaja con cada
acción registrada (trazabilidad).

% ---------------------------------------------------------------- 4
\section{Modelo de datos}

Una sola base SQLite guarda todos los procesos de la planta.

\begin{tabularx}{\linewidth}{@{}P{4.2cm}L@{}}
\toprule
\textbf{Tabla} & \textbf{Contenido}\\
\midrule
\texttt{procesos} & Proceso o línea productiva (nombre, planta, línea, descripción).\\
\texttt{variables} & Característica de calidad y su configuración: tipo de carta, tamaño de muestra $n$, frecuencia, especificaciones (LEI, objetivo, LES), ancho de límites $k$, estimador de $\sigma$, $\sigma$ fijada por el ingeniero, reglas activas, fase, decimales y versión.\\
\texttt{subgrupos} & Cada muestra: número, fecha y hora, fase (\texttt{I}, \texttt{II}, \texttt{archivo}, \texttt{anulado}), conteo y tamaño (atributos), operario, lote, origen (manual, importado, plantilla, semilla) y exclusión con causa y acción.\\
\texttt{mediciones} & Valores individuales de cada subgrupo (variables).\\
\texttt{limites} & Límites congelados para la Fase~II, con historial (uno activo).\\
\texttt{alarmas} & Violaciones de reglas y eventos manuales (parada, retención, anomalía), con explicación en palabras, severidad, estado, causa y acción.\\
\texttt{cambios\_configuracion} & Historial de auditoría: parámetro, valor anterior y nuevo, usuario, justificación, versión y SHA-256 de la configuración.\\
\texttt{lotes\_inspeccion} & Lotes evaluados por muestreo de aceptación y su dictamen.\\
\texttt{reportes} & Reportes generados con su huella SHA-256.\\
\texttt{ajustes} & Ajustes globales (horario de turnos, modo demostración, identificadores de plantillas).\\
\texttt{turnos} & Turnos reales: número, nombre, inicio y fin programados y reales, cómo se cerró (a su hora o a mano) y quién lo abrió y lo cerró; cada subgrupo apunta a su turno.\\
\bottomrule
\end{tabularx}

Las columnas agregadas en versiones posteriores se crean automáticamente al abrir una base
anterior (migración incorporada).

% ---------------------------------------------------------------- 5
\section{Módulos funcionales}

Los módulos 1 a 8 corresponden al alcance técnico pedido para el trabajo. Para cada uno se
indican los métodos implementados.

\subsection{Módulo 1 — Gestión de datos}
\begin{itemize}
\item \textbf{Plantilla de carga paramétrica (recomendada).} La aplicación genera un libro Excel
  a la medida del proceso: una hoja por variable, muestras prenumeradas en dos bloques («para
  establecer los límites» = Fase~I, «monitoreo» = Fase~II), fecha y hora sugeridas según la
  frecuencia, columnas según el tipo de dato, columnas calculadas (promedio, rango,
  \% defectuosas), validación de datos en Excel (solo números, conteos enteros, defectuosas
  $\le$ revisadas), formato condicional para valores fuera de especificación, protección de la
  estructura y una descripción oculta que permite leerla sin mapeos. Al subirla se muestra un
  informe de revisión (errores y avisos por hoja y muestra) antes de guardar; las cargas
  repetidas no duplican muestras y cada muestra entra en su fase.
\item \textbf{Importación de CSV/Excel genéricos:} detección de delimitador y codificación,
  formato largo (una fila por medición) o ancho (una fila por subgrupo), perfilado de columnas,
  mapeo sugerido automáticamente y conversión numérica con coma o punto decimal.
\item \textbf{Limpieza:} valores faltantes (descartar, imputar con la media del subgrupo o
  bloquear la importación) y atípicos (criterio de Tukey con rango intercuartílico, criterio de
  $\pm k\sigma$ o ninguno), con informe de lo modificado.
\item \textbf{Subgrupos racionales:} formación por columna de subgrupo o por tamaño $n$
  consecutivo.
\item \textbf{Tablas resumen:} media, mediana, desviación estándar, mínimo, máximo, cuartiles y
  estadísticos por subgrupo; exportación del conjunto depurado.
\item \textbf{Variables a controlar:} alta de procesos y variables con un asistente de cuatro
  pasos que recomienda la carta según cómo se registra el dato y el tamaño de muestra, y explica
  el porqué.
\end{itemize}

\subsection{Módulo 2 — Cartas de control por variables}

\begin{tabularx}{\linewidth}{@{}P{3cm}L@{}}
\toprule
\textbf{Carta} & \textbf{Límites ($k=3$) y estimador de $\sigma$}\\
\midrule
$\bar{X}$-R \newline ($2\le n\le 10$) & $\bar{\bar{x}} \pm A_2\bar{R}$; \quad R: $[D_3\bar{R},\ D_4\bar{R}]$; \quad $\hat\sigma=\bar{R}/d_2$\\
$\bar{X}$-S \newline ($n>10$ o $n$ variable) & $\bar{\bar{x}} \pm A_3\bar{s}$; \quad S: $[B_3\bar{s},\ B_4\bar{s}]$; \quad $\hat\sigma=\bar{s}/c_4$. Con $n$ variable, límites por subgrupo y $\bar{s}$ combinada.\\
I-MR ($n=1$) & $\bar{x} \pm k\,\overline{MR}/1{,}128$; \quad MR: $\text{LCS}=D_4(2)\,\overline{MR}$\\
\bottomrule
\end{tabularx}

\begin{itemize}
\item Ancho de límites $k$ configurable (fórmulas generales cuando $k\neq 3$) y \textbf{$\sigma$
  fijada por el ingeniero} (límites con estándar dado):
  \[
  \bar{X}:\ \mu \pm k\frac{\sigma}{\sqrt{n}},\qquad
  R:\ (d_2 \pm k\,d_3)\,\sigma,\qquad
  S:\ \bigl(c_4 \pm k\sqrt{1-c_4^2}\bigr)\,\sigma .
  \]
\item \textbf{Fase~I:} cálculo con los subgrupos incluidos, exclusión de subgrupos con causa
  asignable obligatoria, simulación del efecto de excluir antes de confirmar, recálculo
  iterativo y congelamiento de límites (advierte si quedan señales o hay menos de 20 subgrupos).
\item \textbf{Fase~II:} cada subgrupo nuevo se evalúa contra los límites congelados y genera
  alarmas.
\item Constantes $A_2$, $A_3$, $B_3$, $B_4$, $D_3$, $D_4$, $d_2$, $d_3$ y $c_4$ para $n=2$ a $25$.
\end{itemize}

\subsection{Módulo 3 — Cartas de control por atributos}
\begin{itemize}
\item \textbf{p} (fracción defectuosa, $n$ constante o variable con límites por muestra),
  \textbf{np} (número de defectuosos, $n$ constante), \textbf{c} (defectos por muestra) y
  \textbf{u} (defectos por unidad, tamaño variable).
\item Límites, recortados a $[0,1]$ o a $0$:
  \[
  p:\ \bar{p} \pm k\sqrt{\frac{\bar{p}(1-\bar{p})}{n_i}},\qquad
  c:\ \bar{c} \pm k\sqrt{\bar{c}},\qquad
  u:\ \bar{u} \pm k\sqrt{\frac{\bar{u}}{n_i}} .
  \]
\item \textbf{Selección justificada:} un asistente elige la carta según el tipo de dato (unidad
  buena/mala frente a conteo de fallas) y si el tamaño de muestra es constante, con la
  explicación en palabras.
\item Mismo flujo de Fase~I y Fase~II que en variables.
\end{itemize}

\subsection{Reglas de detección (comunes a ambos tipos de carta)}

Se usan las reglas de Nelson N1--N8 y la variante Western Electric WE4, evaluadas sobre el valor
estandarizado $z=(x-\text{LC})/\sigma$ del estadístico graficado, con dos conjuntos predefinidos
(«Nelson» y «Western Electric»).

\begin{tabularx}{\linewidth}{@{}P{1.4cm}L@{}}
\toprule
\textbf{Regla} & \textbf{Patrón}\\
\midrule
N1 & 1 punto fuera de los límites de control\\
N2 & 9 puntos seguidos del mismo lado de la línea central\\
N3 & 6 puntos seguidos crecientes o decrecientes\\
N4 & 14 puntos seguidos alternando arriba y abajo\\
N5 & 2 de 3 puntos más allá de $2\sigma$ del mismo lado\\
N6 & 4 de 5 puntos más allá de $1\sigma$ del mismo lado\\
N7 & 15 puntos seguidos dentro de $\pm1\sigma$ (estratificación)\\
N8 & 8 puntos seguidos fuera de $\pm1\sigma$ a ambos lados (mezcla)\\
WE4 & 8 puntos seguidos del mismo lado de la línea central\\
\bottomrule
\end{tabularx}

Las cartas de dispersión (R, S, MR) evalúan solo N1. Cada violación genera una alarma con
severidad (crítica, advertencia, alerta) y una explicación en palabras simples para el operario;
una racha no genera alarmas repetidas.

\subsection{Módulo 4 — Validación de supuestos}

\begin{tabularx}{\linewidth}{@{}P{3.6cm}L@{}}
\toprule
\textbf{Supuesto} & \textbf{Métodos}\\
\midrule
Normalidad & Shapiro-Wilk, Anderson-Darling, Lilliefors; gráfico Q-Q con bandas de confianza al 95\,\%; asimetría y curtosis.\\
Independencia & Función de autocorrelación (ACF) con bandas $\pm1{,}96/\sqrt{N}$, prueba de Ljung-Box, estadístico de Durbin-Watson.\\
Aleatoriedad y estabilidad & Prueba de rachas sobre la mediana; tendencia por $\tau$ de Kendall.\\
Transformación & Box-Cox con $\lambda$ estimado e intervalo de confianza.\\
\bottomrule
\end{tabularx}

Se emite un veredicto global de aptitud para cartas de Shewhart. Las causas especiales se
identifican con las reglas de detección y se justifican al excluir subgrupos (causa y acción
obligatorias, registradas en el historial).

\subsection{Módulo 5 — Capacidad del proceso}
\begin{itemize}
\item \textbf{Índices:}
  \[
  C_p=\frac{\text{LES}-\text{LEI}}{6\hat\sigma},\qquad
  C_{pu}=\frac{\text{LES}-\hat\mu}{3\hat\sigma},\qquad
  C_{pl}=\frac{\hat\mu-\text{LEI}}{3\hat\sigma},\qquad
  C_{pk}=\min(C_{pu},C_{pl}).
  \]
  Con especificación unilateral se informa el índice del lado especificado.
\item $\hat\sigma$ dentro de subgrupos ($\bar{R}/d_2$, $\bar{s}/c_4$ o $\overline{MR}/1{,}128$) o
  $\sigma$ fijada por el ingeniero.
\item \textbf{Fracción no conforme esperada} bajo normalidad, por lado (debajo del mínimo y
  encima del máximo), en porcentaje y en piezas por millón; \textbf{fracción observada} en los
  datos, por lado.
\item Comparación con la especificación mediante un histograma con curva normal y líneas de
  mínimo, máximo, ideal y promedio.
\item \textbf{Veredicto:} capaz ($C_{pk}\ge 1{,}33$), marginal ($1{,}00\le C_{pk}<1{,}33$) o no
  capaz ($C_{pk}<1{,}00$), con advertencia si el proceso no está bajo control.
\item Para atributos, la capacidad se expresa como el nivel medio de defectos ($\bar{p}$,
  $\bar{c}$ o $\bar{u}$).
\end{itemize}

\subsection{Módulo 6 — Diseño de cartas y plan de muestreo}
\begin{itemize}
\item \textbf{Carta $\bar{X}$ / I}, para un corrimiento de la media de $\delta$ desviaciones
  estándar:
  \begin{gather*}
  \alpha = 2\Phi(-k),\qquad
  \beta = \Phi(k-\delta\sqrt{n})-\Phi(-k-\delta\sqrt{n}),\\
  \text{ARL}_0=\frac{1}{\alpha},\qquad
  \text{ARL}_1=\frac{1}{1-\beta},\qquad
  \text{ATS}=\text{ARL}\cdot h,
  \end{gather*}
  donde $h$ es el intervalo entre muestras; también se informan la desviación estándar de la
  longitud de corrida y la potencia.
\item \textbf{Cartas de atributos:} ARL exacto con la distribución discreta (binomial para p y
  np, Poisson para c y u), con el nivel habitual de defectos y el nivel que se quiere detectar.
\item \textbf{Curva ARL$_1$ (o ATS$_1$) frente al corrimiento $\delta$} para comparar varios
  tamaños de muestra, tabla de desempeño para $\delta=0$ a $3$ y muestras necesarias para
  detectar con 90\,\% de probabilidad.
\item \textbf{Simulación de Monte Carlo} del ARL, con intervalo de confianza y error relativo
  frente al valor teórico.
\item \textbf{Plan de muestreo aplicado a la variable:} se ajustan $n$, frecuencia, $k$ y
  $\sigma$, y se compara el plan vigente con el propuesto en tiempo para detectar un cambio,
  frecuencia de falsas alarmas y \textbf{carga muestral} (unidades por hora y por turno), con
  explicación en palabras. Al guardar el plan en Fase~II, los límites se recalculan con el centro
  y la $\sigma$ congelados; un cambio de $n$ archiva las muestras de monitoreo tomadas con el
  tamaño anterior.
\item \textbf{Carga muestral objetivo} (cartas $\bar{X}$ e I): a partir de la carga actual
  $r = n/h$ (unidades por hora) y una carga objetivo $r^*$, se comparan dos caminos: aumentar $n$
  ($n = r^* h$) o reducir el intervalo ($h = n/r^*$). Para un cambio de la media de $\Delta$
  unidades de la variable ($\delta = \Delta/\sigma$, en cualquier dirección) se informan
  \[
  1-\beta \;(\text{potencia por muestra}),\qquad 1-\beta^{m}\;(\text{primeras } m \text{ muestras}),\qquad
  1-\beta^{\lfloor T/h \rfloor}\;(\text{primeras } T \text{ horas}),
  \]
  además de $\text{ATS}_1 = h/(1-\beta)$ y $\text{ATS}_0 = h/\alpha$, con curvas de potencia y de
  detección acumulada. Cualquiera de los caminos se lleva al plan propuesto para guardarlo.
\end{itemize}

\subsection{Módulo 7 — Muestreo de aceptación}
\begin{itemize}
\item \textbf{Planes simples $(n,c)$:} probabilidad de aceptación
  $P_a=P(d\le c)$ con modelos binomial, hipergeométrico (lote finito $N$) y Poisson.
\item \textbf{Curva característica de operación (OC), calidad media de salida (AOQ) y su máximo
  (AOQL), inspección total media (ATI):}
  \[
  \text{AOQ}=\frac{P_a\,p\,(N-n)}{N},\qquad
  \text{ATI}=n+(1-P_a)(N-n).
  \]
\item \textbf{Diseño del plan} a partir de AQL, LTPD y los riesgos $\alpha$ y $\beta$, por
  búsqueda exacta de $n$ y $c$.
\item \textbf{Dictamen de un lote:} aceptar si $d\le c$, con la fracción defectuosa estimada y su
  intervalo de confianza de Wilson; registro histórico de lotes inspeccionados.
\end{itemize}

\subsection{Módulo 8 — Reportes}
\begin{itemize}
\item \textbf{Informe PDF compilado con LaTeX} (fuente predeterminada de LaTeX), escrito para
  lectores no especialistas y \textbf{solo descriptivo} (no incluye recomendaciones): estado
  general en una frase, hallazgos, indicadores con su significado, carta de control con guía de
  lectura, muestras fuera de lo normal, cumplimiento de la especificación, verificación de
  supuestos, plan de muestreo, alertas, lotes, historial de cambios, anexo técnico y glosario.
  Cuatro plantillas (integral, capacidad, control, trazabilidad) y selección de secciones.
\item \textbf{Formato de respaldo:} si el equipo no tiene LaTeX, el PDF se genera con ReportLab
  y la pantalla lo informa.
\item \textbf{Libro Excel} con una hoja por sección y los datos numéricos reales (subgrupos,
  límites, capacidad, supuestos, diseño, lotes, historial y gráficos).
\item Cada archivo se guarda con su \textbf{huella SHA-256} y puede verificarse después (detecta
  cualquier modificación).
\end{itemize}

\subsection{Componentes transversales}
\begin{itemize}
\item \textbf{Vista Operario:} captura con resultado en vivo (promedio y si está dentro de lo
  normal), aviso de valor probablemente mal escrito (a más de $5\sigma$ del centro habitual o muy
  lejos de la especificación), \textbf{deshacer} el último registro durante 10 minutos (queda como
  «anulado», con motivo), cuenta regresiva a la próxima muestra según la frecuencia (solo dentro
  de un turno), botón para iniciar o terminar el turno, resumen del turno y reporte manual de
  problemas.
\item \textbf{Turnos de trabajo:} el ingeniero define el horario de la planta (nombre, inicio,
  fin, días de trabajo y minutos hasta la primera muestra). Cada turno se abre y se cierra solo a
  sus horas, o a mano desde Operario o Ingeniería, y queda numerado (\#1, \#2\ldots). La próxima
  muestra vence a la última muestra del turno $+\,h$, o al inicio $+$ arranque ($h$ por omisión);
  fuera de turno no hay cuenta regresiva ni muestras atrasadas. Registrar fuera de turno pide
  confirmación y la muestra queda «sin turno». Se rechazan turnos que se superponen el mismo día.
\item \textbf{Alarmas:} seguimiento con estados (sin revisar, en revisión, resuelta); para
  resolver se exigen causa y acción.
\item \textbf{Configuración y auditoría:} especificaciones, reglas, turnos y parámetros con
  versión, motivo obligatorio y SHA-256 de la configuración; historial filtrable y exportable.
\item \textbf{Base de datos y respaldos:} respaldo manual y automático, restauración a un
  respaldo (con respaldo previo automático), copia descargable de la base, restauración desde
  archivo y carga de los datos de demostración en una base vacía.
\end{itemize}

% ---------------------------------------------------------------- 6
\section{Interfaz}\label{sec:interfaz}

\begin{tabularx}{\linewidth}{@{}P{2.1cm}P{4.4cm}L@{}}
\toprule
\textbf{Vista} & \textbf{Pantalla} & \textbf{Función principal}\\
\midrule
Operario & Registrar muestra & Estado del proceso, captura, deshacer, gráfica de las últimas 25 muestras, alertas y turno\\
 & Cómo va el proceso & Gráfica ampliada con guía de lectura y muestras para revisar\\
 & Alertas & Seguimiento de alarmas y eventos\\
 & Historial de muestras & Todas las muestras (incluidas archivadas y anuladas), exportable a CSV\\
\midrule
Ingeniería & Variables a controlar & Procesos y variables; asistente de alta\\
 & Plantilla de carga & Generación y carga de la plantilla Excel paramétrica\\
 & Datos y carga CSV & Importación, limpieza y subgrupos\\
 & Gráficos por variables & Cartas en Fase~I y II, depuración y congelamiento\\
 & Gráficos por atributos & Cartas p, np, c y u en Fase~I y II\\
 & Validación de supuestos & Normalidad, independencia, rachas, Box-Cox\\
 & Capacidad & Índices, fracción no conforme, histograma y veredicto\\
 & Plan de muestreo (ARL) & Diseño de la carta y plan aplicado a la variable\\
 & Muestreo de aceptación & Curvas OC, AOQ, ATI, diseño de planes y dictamen de lotes\\
 & Límites y reglas & Configuración, reglas, turnos, historial de auditoría y modo demostración\\
 & Reportes & Vista previa, PDF, Excel e historial con verificación SHA-256\\
 & Base de datos y respaldos & Ubicación, respaldos, restauración y demostración\\
\bottomrule
\end{tabularx}

Convenciones de diseño: paleta del diseño original (Google Stitch), colores de semáforo
reservados para el estado del proceso, números con coma decimal también en los gráficos y
textos sin jerga en la vista Operario.

% ---------------------------------------------------------------- 7
\needspace{16\baselineskip}
\section{API REST}

La API tiene 72 operaciones bajo el prefijo \texttt{/api}, documentadas en \texttt{/docs}:

\begin{center}
\begin{tabular}{@{}lr@{\qquad\qquad}lr@{}}
\toprule
\textbf{Grupo} & \textbf{Operaciones} & \textbf{Grupo} & \textbf{Operaciones}\\
\midrule
Procesos y variables & 9 & Plantilla de carga & 4\\
Configuración & 10 & Cartas de control & 4\\
Sistema y respaldos & 9 & Diseño ARL/ATS & 3\\
Reportes & 6 & Atributos & 2\\
Subgrupos & 5 & Operario & 1\\
Muestreo de aceptación & 5 & Supuestos & 1\\
Alarmas & 4 & Capacidad & 1\\
Datos & 4 & Turnos & 4\\
\bottomrule
\end{tabular}
\end{center}

Las respuestas son JSON; los valores no numéricos (NaN, $\pm\infty$) se convierten en
\texttt{null}. Las entradas se validan con esquemas Pydantic y los errores se devuelven con
mensajes en español que indican cómo corregirlos.

% ---------------------------------------------------------------- 8
\section{Trazabilidad e integridad}
\begin{itemize}
\item Toda modificación de configuración, exclusión de subgrupos, congelamiento o regreso a
  Fase~I, cambio de plan, importación y anulación de registros queda en el historial con usuario,
  fecha, valor anterior y nuevo, motivo, versión y \textbf{SHA-256 de la configuración
  resultante}.
\item Las muestras no se borran: se excluyen (con causa), se archivan (al reiniciar la Fase~I o
  cambiar $n$) o se anulan (error de captura), y conservan su número.
\item Los reportes y las plantillas importadas se identifican por su huella SHA-256.
\end{itemize}

% ---------------------------------------------------------------- 9
\section{Validación y calidad}
\begin{itemize}
\item \textbf{296 pruebas automáticas} (pytest): cada módulo del núcleo tiene su archivo de
  pruebas, además de pruebas de la API con base temporal, de las pantallas y de la generación de
  archivos.
\item \textbf{Valores de referencia de Montgomery:} cartas $\bar{X}$-R (anillos de pistón), I-MR
  (préstamos), p (jugo de naranja), c (circuitos), u (computadoras), capacidad, ARL y muestreo de
  aceptación. Cuando el libro redondea valores intermedios, la prueba verifica el valor exacto
  con tolerancia y reproduce el redondeo del libro.
\item \textbf{Verificación visual} de las pantallas con Microsoft Edge en modo automático, de la
  plantilla Excel abriéndola en Excel y de los PDF página por página.
\item \textbf{Programa empaquetado} probado de extremo a extremo: instalación, arranque,
  pantallas, descargas, informe con LaTeX y desinstalación.
\end{itemize}

% ---------------------------------------------------------------- 10
\section{Requisitos}

\begin{tabularx}{\linewidth}{@{}P{5cm}L@{}}
\toprule
\textbf{Uso} & \textbf{Requisitos}\\
\midrule
Programa de escritorio & Windows 10 u 11 de 64 bits; Microsoft Edge WebView2 (incluido en Windows~11); $\approx$ 205~MB de disco.\\
Informe PDF con LaTeX (opcional) & MiKTeX (Windows) o TeX Live; sin él se usa el formato de respaldo.\\
Modo servidor / desarrollo & Python 3.11 o superior; $\approx$ 400~MB para el entorno virtual.\\
Construcción del instalador & \texttt{construir.bat} con PyInstaller e Inno Setup 6.\\
\bottomrule
\end{tabularx}

% ---------------------------------------------------------------- 11
\section{Limitaciones y decisiones de alcance}
\begin{itemize}
\item \textbf{Un solo equipo:} el servidor escucha solo en \texttt{127.0.0.1}; no hay acceso
  desde otros equipos de la red ni inicio de sesión (el usuario se identifica por nombre).
\item \textbf{Capacidad:} solo índices $C_p$, $C_{pk}$, $C_{pu}$ y $C_{pl}$ (no se calculan
  $P_p$ ni $P_{pk}$); la fracción no conforme esperada supone distribución normal (por eso se
  informa también la observada y se valida la normalidad).
\item \textbf{Cartas:} Shewhart por variables y atributos; no se incluyen CUSUM ni EWMA.
\item \textbf{Fuera de alcance:} letra código y niveles de inspección de ISO 2859, planes dobles
  y múltiples, ajustes de distribuciones no normales (Johnson, Weibull) e índices no
  paramétricos, firma digital e integración con balanzas o sistemas SCADA.
\item \textbf{Distribución:} el instalador no está firmado digitalmente, por lo que Windows
  muestra una advertencia la primera vez.
\end{itemize}

% ---------------------------------------------------------------- 12
\section{Organización del código}

{\small
\begin{verbatim}
app/
  core/              nucleo estadistico puro
    constantes.py    constantes de cartas (n = 2...25)
    datos.py         lectura, limpieza y subgrupos
    plantilla.py     plantilla Excel parametrica
    graficos_variables.py, graficos_atributos.py, reglas.py
    supuestos.py, capacidad.py, diseno_arl.py, muestreo_aceptacion.py
    reportes.py, reporte_latex.py, figuras.py, plantillas/reporte.tex
  api/               FastAPI: un router por modulo + servicios.py
  db/                modelos, sesion, migraciones, demostracion, respaldos
  web/               pantallas, plantillas Jinja2, JavaScript, recursos locales
  escritorio.py      programa de escritorio (ventana + servidor local)
  iniciar.py         arranque en modo servidor
instalador/          receta de PyInstaller, script de Inno Setup e icono
tests/               pruebas automaticas
datos_prueba/        conjuntos de datos de ejemplo y generador
docs/                instalacion, manual, esta descripcion y capturas
\end{verbatim}
}

\end{document}
